\chapter{Background}
\label{ch:background}

This chapter establishes the technical foundations on which the remainder of
this book depends. Section~\ref{sec:bg-pm} introduces process mining:
its vocabulary, formal models, and the three-phase lifecycle of discovery,
conformance, and enhancement. Section~\ref{sec:bg-ocpm} motivates
object-centric process mining and presents the OCEL~2.0 formal model.
Section~\ref{sec:bg-rust} explains the Rust type-system features that make
compile-time law encoding feasible. Section~\ref{sec:bg-verification} surveys
the landscape of type-based verification and locates nightly Rust within it.
Section~\ref{sec:bg-trybuild} describes the trybuild harness and the
type-law receipt protocol it enables.

% -----------------------------------------------------------------------------
\section{Process Mining}
\label{sec:bg-pm}
% -----------------------------------------------------------------------------

Process mining is a discipline at the intersection of data science and process
management that aims to discover, monitor, and improve real processes by
extracting knowledge from \emph{event logs}~\cite{vanderaalst2016process}.
Where classical business-process management assumes a prescribed model and
checks whether the organisation follows it, process mining starts from
observed behaviour --- timestamps, activity labels, case identifiers --- and
derives models, deviations, and improvement opportunities from first
principles.

\paragraph{Event logs and traces.}
An \emph{event} is the smallest observable unit: it records that a specific
activity occurred at a specific time, optionally linked to a case, resource,
and payload of typed attributes. A \emph{trace} is a finite, time-ordered
sequence of events sharing a common case identifier. An \emph{event log} is a
collection of traces. This vocabulary --- event, trace, log --- originated with the
\textsc{WorkflowNet} formalism~\cite{vanderaalst1998application} and was
standardised in the eXtensible Event Stream (\textsc{XES}) format by the
\textsc{IEEE} Task Force on Process Mining~\cite{gunther2014xes}.

\begin{definition}[Event Log]
\label{def:event-log}
Let $\mathcal{A}$ be a universe of \emph{activity labels} and $\mathcal{T}$ a
totally ordered universe of \emph{timestamps}. An \emph{event} is a tuple
$e = (a, t, \mathbf{d})$ where $a \in \mathcal{A}$, $t \in \mathcal{T}$, and
$\mathbf{d}$ is a finite map of typed attribute data. A \emph{trace} over
$\mathcal{A}$ is a finite sequence $\sigma = \langle e_1, e_2, \ldots, e_n
\rangle$ of events with $t(e_1) \leq t(e_2) \leq \cdots \leq t(e_n)$.
An \emph{event log} $L$ over a universe of events $\mathcal{E}$ is a multiset
$L \in \mathcal{B}(\mathcal{A}^*)$ such that every element of $L$ is a trace
over $\mathcal{A}$.
\end{definition}

The multiset formulation is important: a log may contain duplicate traces.
This occurs whenever the same sequence of activities is observed for multiple
cases, and collapsing duplicates would misrepresent the frequency distribution
of process variants --- a key input to performance analysis.

\paragraph{The PM lifecycle.}
Van der Aalst identifies three phases of the process-mining
lifecycle~\cite{vanderaalst2016process}:
\begin{enumerate}
  \item \textbf{Discovery.} An algorithm takes a log $L$ as input and produces
        a process model $M$ that describes the observed behaviour.
        The Alpha Miner~\cite{vanderaalst2004workflow}, the Inductive
        Miner~\cite{leemans2013discovering}, and the Split
        Miner~\cite{augusto2019split} are representative algorithms, each
        with different guarantees on soundness, fitness, and precision.
  \item \textbf{Conformance checking.} Given both a model $M$ and a log $L$,
        conformance checking measures how well $M$ and $L$ agree.
        Alignment-based techniques~\cite{adriansyah2011conformance} compute
        an optimal alignment between each trace and the model,
        assigning cost to deviations and producing per-event diagnostics.
        The four quality dimensions --- fitness, precision, generalisation, and
        simplicity~\cite{buijs2014quality} --- are the canonical scorecard.
  \item \textbf{Enhancement.} Using both the model and the log, enhancement
        extends or repairs the model: annotating it with performance data,
        repairing deviating fragments, or updating it to reflect observed
        drift~\cite{vanderaalst2016process}.
\end{enumerate}

\paragraph{XES vs OCEL~2.0.}
The XES standard~\cite{gunther2014xes} encodes case-centric event logs in XML.
Every event belongs to exactly one trace identified by a single case attribute.
XES is well-supported by toolkits such as pm4py~\cite{berti2019pm4py} and
ProM~\cite{verbeek2010prom}, and remains the dominant interchange format for
classical process mining. Its central limitation is the \emph{single case
notion}: real processes involve multiple interacting objects (orders, items,
deliveries, payments) and forcing all events into a single case dimension
requires an arbitrary choice of which object to use as the case, discarding
relational structure in the process.

The Object-Centric Event Log format (OCEL~2.0,~\cite{ocel20_2023}) replaces
the case notion with a typed object graph. Each event may relate to any number
of objects of any type, and objects may relate to each other.
Section~\ref{sec:bg-ocpm} formalises this model.

% -----------------------------------------------------------------------------
\section{Object-Centric Process Mining}
\label{sec:bg-ocpm}
% -----------------------------------------------------------------------------

Classical process mining assumes that every event belongs to exactly one case.
This assumption is adequate when a process is genuinely case-centric ---
insurance claims, for example, follow a single claimant through a sequence of
activities and a single case identifier captures the trajectory cleanly.
But most real processes involve \emph{convergence} and \emph{divergence}: a
sales order converges onto a shipment that contains multiple items, and a
single payment diverges to satisfy multiple invoices. Flattening such a process
to a single case notion requires choosing one object type as the case ---
say, the order --- and then either duplicating events (once per item in the
order) or dropping them (once per order, ignoring per-item variation).
Both choices introduce artefacts: duplication inflates frequency counts;
dropping loses causal detail~\cite{ghahfarokhi2021ocel}.

Object-centric process mining (OCPM) resolves this by treating the event graph
as a first-class citizen: events relate to multiple objects across multiple
types, and objects relate to each other, without any flattening~\cite{aalst2020ocpm}.
The formal model is OCEL~2.0.

\begin{definition}[OCEL~2.0 Event Log]
\label{def:ocel20}
Let $\mathcal{OT}$ be a finite universe of \emph{object types} and
$\mathcal{A}$ a universe of \emph{activity labels}. An \emph{OCEL~2.0 event
log} is a tuple $\mathcal{O} = (E, O, \mathit{EA}, \mathit{OA}, E_2O, O_2O)$
where:
\begin{itemize}
  \item $E$ is a finite set of \emph{events}, each bearing an activity label
        $a \in \mathcal{A}$ and a timestamp $t \in \mathcal{T}$.
  \item $O$ is a finite set of \emph{objects}, each bearing a type
        $\tau \in \mathcal{OT}$ and a typed attribute map.
  \item $\mathit{EA} : E \to (\mathit{Key} \rightharpoonup \mathit{Val})$ maps
        events to typed attribute records.
  \item $\mathit{OA} : O \times \mathcal{T} \to (\mathit{Key}
        \rightharpoonup \mathit{Val})$ maps objects and timestamps to attribute
        records, capturing \emph{object change} over time.
  \item $E_2O \subseteq E \times O \times \mathit{Qual}$ is the
        \emph{event-to-object relation} (E2O), where $\mathit{Qual}$ is a
        finite universe of \emph{qualifiers} naming the role of an object in
        an event (e.g. \texttt{"item"}, \texttt{"responsible"}).
  \item $O_2O \subseteq O \times O \times \mathit{Qual}$ is the
        \emph{object-to-object relation} (O2O), capturing typed structural
        relationships between objects (e.g. \texttt{"part\_of"},
        \texttt{"follows"}).
\end{itemize}
The log is \emph{well-formed} if every $(e, o, q) \in E_2O$ satisfies
$e \in E$ and $o \in O$, and every $(o_1, o_2, q) \in O_2O$ satisfies
$o_1, o_2 \in O$.
\end{definition}

The qualification of relations --- the $\mathit{Qual}$ dimension on both $E_2O$
and $O_2O$ --- distinguishes OCEL~2.0 from earlier object-centric proposals.
Earlier formats~\cite{ghahfarokhi2021ocel} recorded only \emph{which} objects
an event related to, not \emph{in what role}. Qualifiers enable richer
downstream analysis: an alignment algorithm can distinguish events where an
order was the \emph{target} of a payment from events where it was merely a
\emph{reference}, without relying on event-attribute heuristics.

\paragraph{The \texttt{wasm4pm-compat} representation.}
The crate's \texttt{ocel} module gives OCEL its own genuine canon rather than
modelling it as ``an \texttt{EventLog} with side tables''.
\texttt{EventObjectLink} (E2O) and \texttt{ObjectObjectLink} (O2O) are
first-class types in the module. The \texttt{OcelLog::validate} method performs
structural integrity: every link must reference declared objects and events;
identifiers must be unique. Validation is structural --- it does not discover
object-centric Petri nets, compute OC-DFGs, or check conformance.
Those functions graduate to the \texttt{wasm4pm} execution engine.

The flattening trap is a named law, \texttt{OcelRefusal::FlatteningLoss}:
any projection from OCEL to a case-centric shape must carry a
\texttt{LossPolicy} and emit a \texttt{LossReport}. Silent structure loss
is a compile-time defect, not a documentation note.

% -----------------------------------------------------------------------------
\section{The Rust Type System}
\label{sec:bg-rust}
% -----------------------------------------------------------------------------

Rust is a systems programming language whose type system provides memory safety
without a garbage collector, enforced by the \emph{ownership} and
\emph{borrow checker} at compile time~\cite{matsakis2014rust}. For this work,
three features of the type system are foundational: the trait system, const
generics, and \texttt{PhantomData}. Three nightly-only extensions elevate these
features from general-purpose machinery to a type-law enforcement substrate.

\paragraph{Ownership and the borrow checker.}
Every value in Rust has exactly one owner. Ownership may be transferred
(\emph{move}) or temporarily lent (\emph{borrow}), but the borrow checker
statically prevents data races and use-after-free. For process evidence, the
ownership model has an important consequence: a value of type
\texttt{Evidence<T, Raw, W>} cannot be ``shared'' into an
\texttt{Evidence<T, Admitted, W>} context by accident; the types are nominal
and distinct, and the borrow checker will not silently coerce one to the other.

\paragraph{The trait system.}
Traits are Rust's mechanism for bounded polymorphism: they name a set of
associated types and method signatures, and a type implements a trait by
providing concrete definitions for each. Traits support \emph{sealed} patterns:
a trait defined in crate $C$ with no public implementation path outside $C$
cannot be implemented by downstream crates. \texttt{wasm4pm-compat} uses sealed
traits for the \texttt{Witness} family and for the \texttt{Admit} lifecycle
gate, ensuring that only crate-internal paths can produce admitted evidence.

\paragraph{Const generics and \texttt{PhantomData}.}
Const generics allow integer, boolean, or enum values to appear as type
parameters. A type \texttt{WfNetConst<SOUNDNESS: bool>} is a distinct type for
each boolean value of \texttt{SOUNDNESS}. \texttt{PhantomData<T>} is a
zero-sized field that makes a struct \emph{generic over} \texttt{T} without
storing any \texttt{T} data; the distinction is entirely at the type level and
is erased at compile time. State tokens (\texttt{Raw}, \texttt{Admitted},
\texttt{Refused}, \ldots) and witness markers (\texttt{Ocel20},
\texttt{Xes1849}, \ldots) are both encoded as \texttt{PhantomData}: they
contribute zero bytes to the memory layout of \texttt{Evidence<T, S, W>}
while remaining fully opaque to the type checker.

\paragraph{Nightly features.}
The crate requires four nightly features that are currently undergoing
stabilisation:

\begin{itemize}
  \item \texttt{generic\_const\_exprs} --- allows boolean expressions over const
        generic parameters in where-bounds. Without this feature, a bound such
        as \texttt{where Require<\{ N <= 8 \}>: IsTrue} cannot be expressed.
  \item \texttt{adt\_const\_params} --- allows arbitrary \texttt{ConstParamTy}
        enum values as const generic parameters, not just integers and booleans.
        This enables \texttt{WfNetConst<SOUNDNESS>} where \texttt{SOUNDNESS} is
        an enum variant naming a specific soundness criterion.
  \item \texttt{const\_trait\_impl} --- allows trait methods to be evaluated in
        const contexts. This is required for the \texttt{Between01} bounds that
        enforce metric values at const-evaluation time.
  \item \texttt{min\_specialization} --- allows a limited form of trait
        specialisation: a more specific impl may override a more general one.
        Used in the conformance module to provide specialised behaviour for
        metrics known at type level to be exact.
\end{itemize}

\paragraph{The Assert/IsTrue pattern.}
The compile-time law kernel at the heart of \texttt{wasm4pm-compat} is the
\texttt{Assert}/\texttt{IsTrue}/\texttt{Require} trio, shown below:

\begin{lstlisting}[language=Rust, caption={Assert/IsTrue compile-time law machinery from
  \texttt{src/law.rs}. \texttt{Assert<true>} is inhabited; \texttt{Assert<false>}
  is not. A where-bound \texttt{Require<\{EXPR\}>: IsTrue} becomes a compile
  error when \texttt{EXPR} evaluates to \texttt{false}.},
  label={lst:bg-assert-istrue}]
pub struct Assert<const OK: bool>;
pub trait IsTrue {}
impl IsTrue for Assert<true> {}
pub type Require<const OK: bool> = Assert<OK>;
\end{lstlisting}

The mechanism is as follows. \texttt{Assert<true>} is an inhabited struct
(it can be constructed as \texttt{Assert}). \texttt{Assert<false>} is also
a valid type, but \texttt{IsTrue} is only implemented for \texttt{Assert<true>}.
Therefore, any function \texttt{fn f<const N: usize>() where Require<\{ N > 0
\}>: IsTrue} can only be called when $N > 0$; calling it with $N = 0$ produces
a compile error because \texttt{Assert<false>: IsTrue} has no proof. This
mechanism extends to arbitrarily complex const expressions:
\texttt{ConditionCell<BITS>} uses \texttt{Require<\{ BITS <= 8 \}>: IsTrue}
to enforce the \emph{Need9 means split} law from the Blue River Dam covenant,
and \texttt{Between01<NUM, DEN>} uses \texttt{Require<\{ NUM <= DEN \}>: IsTrue}
to enforce metric-range soundness. In every case, the enforcement is total and
zero-cost: it occurs at monomorphisation time and leaves no trace at runtime.

% -----------------------------------------------------------------------------
\section{Verification in Type Systems}
\label{sec:bg-verification}
% -----------------------------------------------------------------------------

The idea that programs should carry machine-checkable proofs of their own
correctness has a long history. The goal of this section is not to survey that
history exhaustively but to locate the \texttt{wasm4pm-compat} approach
precisely within it --- explaining what it shares with and what it differs from
the major traditions.

\paragraph{Dependent types.}
Dependently typed languages such as Coq~\cite{coq2022}, Agda~\cite{norell2009agda},
and Idris~\cite{brady2013idris} allow types to depend on values: the type
\texttt{Vec A n} is the type of vectors over \texttt{A} of \emph{length
exactly} $n$, where $n$ is a runtime value that has been lifted into the type.
This expressive power enables specifications that are virtually unlimited in
complexity: a sorting function can be typed so that its output is not merely a
list but a permutation of its input that is monotone with respect to a provided
order. The cost is proof burden: every dependently typed program is simultaneously
a proof, and the programmer must explicitly construct the proof terms. For most
software engineering contexts, this overhead is prohibitive.

\paragraph{Liquid types.}
Liquid types~\cite{rondon2008liquid} (and their descendant, Liquid
Haskell~\cite{vazou2014liquidhaskell}) add refinement predicates to types
without requiring full dependent types. A liquid type \texttt{\{x : Int |
x > 0\}} is the type of positive integers. The type checker discharges
refinement obligations to an SMT solver, reducing the proof burden to writing
refinement annotations rather than proof terms. Liquid types are more
expressive than simple Hindley--Milner types and less expressive than full
dependent types; they occupy a practical middle ground for numeric and
container invariants.

\paragraph{GADTs.}
Generalised algebraic data types (GADTs) in Haskell~\cite{peytonjones2006gadts}
allow constructors to return specialised instantiations of a parameterised
type. They enable a form of typestate encoding: a value of type
\texttt{Evidence Raw} is distinct from \texttt{Evidence Admitted}, and pattern
matching on the constructor reveals which state the value is in. GADTs are
expressive enough to encode lifecycle invariants but do not support the
const-numeric properties (``this metric is between 0 and 1'',
``this condition cell has at most 8 bits'') that \texttt{wasm4pm-compat} needs.

\paragraph{The gap in process mining implementations.}
Despite this rich landscape, most process mining implementations offer no
compile-time verification of their structural invariants.
pm4py~\cite{berti2019pm4py} represents a Petri net as a mutable Python object;
nothing prevents a caller from constructing a net without a final marking,
naming it sound, and passing it to an alignment algorithm that assumes
soundness. The OCEL~2.0 specification~\cite{ocel20_2023} defines precise
well-formedness conditions, but a JSON deserialiser that produces an OCEL
object has no obligation to check them before returning. Runtime assertions
can catch violations in test runs; they cannot guarantee that violations
cannot occur in code paths that tests do not cover.

\paragraph{How nightly Rust fills the gap.}
Nightly Rust is not a dependently typed language, and it is not a liquid-type
system. It cannot express arbitrary refinement predicates or term-level proofs.
What it can express is a well-defined subset of structural invariants that are
decidable at monomorphisation time: boolean combinations of const generic
inequalities, closed-world trait bounds, sealed lifecycle gates, and phantom
witness types. This subset is precisely what process-mining type law requires.
The invariant ``a log cannot be used as admitted before being admitted against
a named standard'' is a trait-bound invariant, not an arbitrary predicate, and
nightly Rust enforces it completely. The invariant ``a conformance metric is in
$[0,1]$'' is a const-generic inequality, and \texttt{Between01<NUM, DEN>}
encodes it completely. Together, these mechanisms cover the structural core of
OCEL~2.0, WF-net soundness, POWL tree arity, condition-cell capacity, and the
evidence lifecycle --- all at zero runtime cost, all checked by the compiler.

% -----------------------------------------------------------------------------
\section{trybuild and Type-Law Receipts}
\label{sec:bg-trybuild}
% -----------------------------------------------------------------------------

Unit tests and integration tests verify that a program produces the correct
output for a given input. They cannot directly verify that a program \emph{does
not compile} when an invariant is violated --- because by definition, a program
that does not compile has no output to check. The trybuild
crate~\cite{tolnay2019trybuild} fills this gap: it is a test harness for
compile-time behaviour.

\paragraph{How trybuild works.}
A trybuild test suite consists of a collection of small Rust source files,
each of which is compiled independently by the trybuild runner. For each file,
the runner records whether compilation succeeded or failed, and if it failed,
what the compiler diagnostic was. The suite declares a set of files as
\emph{compile-fail fixtures} and a set as \emph{compile-pass fixtures}.
A compile-fail fixture is expected to fail; if it succeeds, the test fails.
A compile-pass fixture is expected to succeed; if it fails, the test fails.
For compile-fail fixtures, trybuild compares the actual diagnostic against a
\emph{golden file} (a \texttt{.stderr} snapshot) stored alongside the fixture.
If the diagnostic changes --- because the compiler produced a different error
message --- the test fails, and the fixture must be updated before the suite
passes.

\paragraph{The .stderr receipt protocol.}
The \texttt{.stderr} golden file is not merely a regression guard.
In the \texttt{wasm4pm-compat} manufacturing doctrine, it is a
\emph{type-law receipt}: a sealed, machine-readable attestation that the
compiler rejected a specific violation of a named law. The key property is
specificity. A compile-fail fixture that fails because of a typo or a missing
import is not a valid receipt: it does not demonstrate that the law is enforced,
only that the program is malformed. The \texttt{.stderr} file pins the exact
error code, span, and message, ensuring that if the law \emph{stops} being
enforced --- because a type bound was accidentally weakened in a refactor --- the
golden file no longer matches the actual diagnostic, and the test fails.

\paragraph{Why compile-fail fixtures are stronger evidence than unit tests.}
A unit test for a type-law invariant is necessarily indirect: it must reach
the invariant through runtime code paths, which means it only tests the paths
the test author thought to exercise. A compile-fail fixture is \emph{total}:
if the program in the fixture compiles without error, the invariant has been
violated for every possible use of the guarded type, not just the paths
the test exercises. Compile-fail fixtures also survive dead-code elimination
and test-coverage gaps because they test the type system, not the runtime.

\begin{definition}[Type-Law Receipt]
\label{def:type-law-receipt}
A \emph{type-law receipt} $R$ for law $\mathcal{L}$ is a pair
$(F^-, \sigma)$ where $F^-$ is a compile-fail fixture that encodes a
violation of $\mathcal{L}$, and $\sigma$ is the sealed compiler diagnostic
confirming the rejection, such that:
\begin{enumerate}
  \item $F^-$ compiles to a type error (not a syntax error, missing-import
        error, or unrelated diagnostic);
  \item $\sigma$ names $\mathcal{L}$ by its Rust error code and span, and
        the expected message is a substring of the actual compiler output;
  \item there exists a corresponding compile-pass fixture $F^+$ that encodes
        the lawful path and compiles without error, proving that $\mathcal{L}$
        does not over-constrain the type.
\end{enumerate}
A receipt is \emph{valid} if and only if all three conditions hold simultaneously.
A receipt that satisfies only (1) and (2) but lacks a corresponding $F^+$ is
a \emph{partial receipt}: it proves rejection but not satisfiability.
\end{definition}

\paragraph{The ALIVE gate.}
The \texttt{wasm4pm-compat} ALIVE gate is the collection of all valid receipts
in \texttt{tests/ui/}: compile-fail fixtures in \texttt{tests/ui/compile\_fail/}
paired with \texttt{.stderr} golden files, and compile-pass fixtures in
\texttt{tests/ui/compile\_pass/}. The gate is run via:

\begin{lstlisting}[language=Rust, caption={Running the ALIVE gate. The
  \texttt{--ignored} flag is required because trybuild fixtures are marked
  \texttt{\#[ignore]} to exclude them from the fast dev-loop test run.},
  label={lst:bg-alive-gate}]
cargo test --test ui_tests -- --ignored
\end{lstlisting}

The gate is deliberately excluded from the default \texttt{cargo test} run
because each fixture is an independent \texttt{rustc} invocation and the
full suite of 602 fixtures takes several minutes to complete. It is the
\emph{certification} surface, not the \emph{development} surface. Chapter
\ref{ch:alive} describes the gate, its fixture disciplines, and the four
certification levels (\textsc{Alive 001} through \textsc{Crown Alive 004})
in detail.
